\chapter{Diseño de la solución}
\label{chap:diseno}

\ph{Este capítulo describe \textbf{cómo} se ha diseñado el sistema para cumplir
los requisitos del capítulo anterior. Se centra en la arquitectura, los modelos
y los flujos de datos, sin entrar todavía en el código concreto (eso es la
implementación).}

\section{\ph{Modelo conceptual del dominio}}
\label{sec:modelo_dominio}

\ph{Si el TFG se apoya en un modelo formal del dominio (modelo de features,
ontología, modelo de datos, gramática, modelo matemático\dots), descríbelo aquí.
Si no aplica, sustituye esta sección por la que mejor encaje con el trabajo.}

\ph{Cubre:}
\begin{itemize}
  \item \ph{Para qué sirve el modelo y qué representa.}
  \item \ph{Dónde se almacena (archivo, formato, lenguaje).}
  \item \ph{Cómo se relaciona con el resto del sistema (validación, generación, persistencia).}
\end{itemize}

\subsection{\ph{Estructura del modelo}}
\label{subsec:estructura_modelo}

\ph{Describe la jerarquía o estructura principal del modelo. Si hay un árbol,
descríbelo de raíz a hojas. Si hay entidades, descríbelas con sus relaciones.}

\subsection{\ph{Restricciones del dominio}}
\label{subsec:constraints_dominio}

\ph{Restricciones que el modelo no captura por estructura y que se expresan
como reglas adicionales (rangos, dependencias, exclusiones, fórmulas\dots).
Justifica brevemente cada una desde el punto de vista del dominio.}

\begin{itemize}
  \item \textbf{\ph{Restricción 1}:} \ph{descripción y justificación.}
  \item \textbf{\ph{Restricción 2}:} \ph{descripción y justificación.}
  \item \textbf{\ph{Restricción 3}:} \ph{descripción y justificación.}
\end{itemize}

\subsection{\ph{Espacio de soluciones / configuraciones}}
\label{subsec:espacio_soluciones}

\ph{Si tiene sentido, calcula o estima el tamaño del espacio de soluciones
válidas (combinatoria de opciones, rango continuo de parámetros, etc.).
Es útil para ilustrar la magnitud del problema.}

\section{Arquitectura del sistema}
\label{sec:arquitectura}

\subsection{Visión general}
\label{subsec:arquitectura_general}

\ph{Describe la arquitectura de alto nivel: tipo de aplicación
(monolito, microservicios, librería, CLI\dots), patrón principal
(capas, hexagonal, MVC\dots) y las tecnologías centrales.}

\ph{Lista de capas habituales:}

\begin{itemize}
  \item \textbf{Capa de presentación}: \ph{interfaz visible al usuario.}
  \item \textbf{Capa de API / controladores}: \ph{coordinación de peticiones y respuestas.}
  \item \textbf{Capa de dominio / lógica}: \ph{reglas de negocio, validación, generación.}
  \item \textbf{Capa de persistencia}: \ph{acceso a base de datos o almacenamiento.}
  \item \textbf{Capa de infraestructura}: \ph{contenedores, despliegue, CI/CD.}
\end{itemize}

\ph{Si usas un modelo de documentación arquitectónica concreto (C4, 4+1, ATAM\dots),
nómbralo aquí y comenta qué niveles incluyes.}

\ph{(Espacio reservado para los diagramas de contexto y de contenedores,
no incluidos en esta plantilla.)}

\subsection{Diagrama de componentes}
\label{subsec:diagrama_componentes}

\ph{Describe cómo se distribuyen los módulos del sistema entre las capas
y cómo se comunican entre sí. Indica responsabilidades por módulo.}

\ph{(Espacio reservado para el diagrama de capas / componentes.)}

\subsection{Diagrama de secuencia del flujo principal}
\label{subsec:diagrama_secuencia}

\ph{Enumera, paso a paso, el flujo principal del sistema (el caso de uso central).
Esto orienta al lector antes de los detalles de implementación.}

\begin{enumerate}
  \item \ph{Paso 1.}
  \item \ph{Paso 2.}
  \item \ph{Paso 3.}
  \item \ph{Paso 4.}
  \item \ph{Paso 5.}
  \item \ph{Paso 6.}
\end{enumerate}

\ph{Comenta también qué ocurre en caso de error (el flujo se interrumpe,
se devuelve un código de error, se notifica al usuario\dots).}

\section{Patrones de diseño aplicados}
\label{sec:patrones}

\ph{Enumera los patrones de diseño (creacionales, estructurales, de
comportamiento o arquitectónicos) que se han aplicado en la solución.
Para cada uno indica \emph{qué problema resuelve}, \emph{dónde se aplica}
en el sistema y \emph{por qué} se eligió frente a alternativas.}

\ph{Ejemplos habituales en TFGs de Ingeniería del Software:}

\begin{itemize}
  \item \textbf{\ph{Patrón 1 (p.\,ej.\ MVC, Repositorio, Singleton, Factory\dots)}:}
        \ph{problema que resuelve, módulo o capa donde se aplica, alternativas
        descartadas.}
  \item \textbf{\ph{Patrón 2}:} \ph{descripción y justificación.}
  \item \textbf{\ph{Patrón 3}:} \ph{descripción y justificación.}
\end{itemize}

\ph{Si la arquitectura sigue un patrón global (capas, hexagonal, MVC,
microservicios\dots), descríbelo aquí o referencia la sección de arquitectura
donde se haya tratado.}

\section{Diseño de la interfaz de usuario}
\label{sec:ui_ux}

\ph{Si el sistema tiene interfaz gráfica (web, móvil, escritorio), dedica una
sección al diseño de la experiencia de usuario. No basta con decir ``se ha
hecho bonito''; conviene documentar las decisiones tomadas.}

\subsection{Mapa de navegación}
\label{subsec:mapa_navegacion}

\ph{Diagrama o lista jerárquica de las pantallas/vistas del sistema y los
caminos posibles entre ellas. Sirve para que el lector se haga una idea de la
extensión funcional sin recorrerse el sistema en vivo.}

\ph{Listado de vistas:}

\begin{itemize}
  \item \ph{Vista 1 — propósito y enlaces salientes.}
  \item \ph{Vista 2 — propósito y enlaces salientes.}
  \item \ph{Vista 3 — propósito y enlaces salientes.}
\end{itemize}

\subsection{Wireframes y mockups}
\label{subsec:mockups}

\ph{Bocetos de baja fidelidad (wireframes) o prototipos visuales (mockups)
de las pantallas principales. Conviene incluir al menos las dos o tres más
representativas: la pantalla principal, la pantalla del flujo principal y
una pantalla de gestión/listado.}

\ph{(Espacio reservado para los wireframes/mockups, no incluidos en esta plantilla.)}

\subsection{Decisiones de diseño visual}
\label{subsec:decisiones_visuales}

\ph{Comenta brevemente las decisiones globales de la interfaz:}

\begin{itemize}
  \item \textbf{\ph{Paleta de colores}}: \ph{principales, secundarios,
        accesibilidad de contraste.}
  \item \textbf{\ph{Tipografía}}: \ph{familias usadas y por qué.}
  \item \textbf{\ph{Sistema de diseño / framework CSS}}: \ph{si se ha usado uno
        (Bootstrap, Tailwind, Material\dots) o se han creado componentes propios.}
  \item \textbf{\ph{Responsive / accesibilidad}}: \ph{breakpoints, criterios de
        accesibilidad seguidos (contraste, navegación por teclado, ARIA\dots).}
\end{itemize}

\section{Esquema de base de datos}
\label{sec:esquema_db}

\subsection{Modelo entidad-relación}
\label{subsec:er}

\ph{Describe las entidades persistidas, sus atributos y las relaciones entre
ellas. Suele bastar con un par de tablas y una representación simbólica de
las relaciones.}

% Hueco para la representación simbólica de las relaciones. Descomenta y adapta:
% \[
% \texttt{entidad\_A} \;\; 1 \longrightarrow N \;\; \texttt{entidad\_B}
% \]

\ph{(Espacio reservado para la representación simbólica de las relaciones,
p.~ej.\ una relación de uno a muchos: \texttt{entidad\_A} $1 \rightarrow N$
\texttt{entidad\_B}. Descomenta y adapta el bloque comentado del fuente.)}

\begin{table}[H]
    \centering
    \begin{tabular}{lll}
        \toprule
        \textbf{Campo} & \textbf{Tipo de dato} & \textbf{Restricciones} \\
        \midrule
        \ph{campo\_1} & \ph{Tipo} & \ph{PK / NOT NULL / \dots} \\
        \ph{campo\_2} & \ph{Tipo} & \ph{\dots} \\
        \ph{campo\_3} & \ph{Tipo} & \ph{\dots} \\
        \ph{campo\_4} & \ph{Tipo} & \ph{\dots} \\
        \bottomrule
    \end{tabular}
    \caption{Esquema de la tabla \ph{\texttt{nombre\_tabla\_1}}.}
    \label{tab:esquema_tabla_1}
\end{table}

\begin{table}[H]
    \centering
    \begin{tabular}{lll}
        \toprule
        \textbf{Campo} & \textbf{Tipo de dato} & \textbf{Restricciones} \\
        \midrule
        \ph{campo\_1} & \ph{Tipo} & \ph{PK / FK / \dots} \\
        \ph{campo\_2} & \ph{Tipo} & \ph{\dots} \\
        \ph{campo\_3} & \ph{Tipo} & \ph{\dots} \\
        \bottomrule
    \end{tabular}
    \caption{Esquema de la tabla \ph{\texttt{nombre\_tabla\_2}}.}
    \label{tab:esquema_tabla_2}
\end{table}

\ph{Comenta cómo se garantiza el aislamiento entre usuarios (filtros por
\texttt{user\_id}, comprobaciones de propiedad, etc.) si aplica.}

\subsection{Migraciones / versionado del esquema}
\label{subsec:migraciones}

\ph{Si el esquema se versiona (Alembic, Flyway, Liquibase\dots), explica
brevemente la herramienta usada, dónde se almacenan las migraciones y qué
comandos se utilizan en desarrollo y en CI.}

\section{Flujo de datos principal}
\label{sec:flujo}

\ph{Describe el flujo de transformación de datos desde la entrada del usuario
hasta el artefacto final. Por ejemplo:
\textbf{Entrada} $\rightarrow$ \textbf{Representación intermedia}
$\rightarrow$ \textbf{Artefacto final}.}

\subsection{\ph{Estructura de la representación intermedia}}
\label{subsec:representacion_intermedia}

\ph{Si el sistema produce una representación intermedia
(JSON canónico, IR, AST, fichero de configuración\dots), describe sus campos
principales y referencia el apéndice donde se incluya un ejemplo completo.}

\subsection{\ph{Mapeo entre niveles}}
\label{subsec:mapeo}

\ph{Si el flujo implica transformaciones explícitas (símbolos a valores
numéricos, alias a constantes, claves a entidades\dots), describe los mapeos
con ejemplos concretos.}

\begin{itemize}
  \item \ph{\texttt{símbolo\_a} $\rightarrow$ \texttt{valor\_a}}
  \item \ph{\texttt{símbolo\_b} $\rightarrow$ \texttt{valor\_b}}
\end{itemize}

\subsection{\ph{Trazabilidad e identificadores}}
\label{subsec:trazabilidad}

\ph{Si el sistema necesita relacionar artefactos generados con su configuración
de origen, explica el esquema de identificadores (UUID, hash, slug\dots)
y dónde se persisten.}

% --- Transición al siguiente capítulo ---
\ph{Cierra el capítulo con un párrafo puente de 2--4 líneas: qué queda
establecido aquí y qué pregunta abre el capítulo siguiente. Por ejemplo:
«Definidos la arquitectura, los modelos y los flujos de datos, el diseño
queda cerrado sobre el papel. El Capítulo~\ref{chap:implementacion} desciende
al código: las tecnologías elegidas y cómo cada módulo materializa las
decisiones tomadas aquí».}
